% TODO make hw stylesheet

\documentclass[10pt]{article}
\usepackage[letterpaper,top=1in,left=1in,right=1in]{geometry}
\usepackage[dvipsnames,svgnames]{xcolor}
\usepackage[shortlabels]{enumitem}
\usepackage[framemethod=TikZ]{mdframed}
\usepackage{amsmath,amssymb,amsthm}
\usepackage[colorlinks]{hyperref}
\usepackage{multicol}
\usepackage{tabularx}
\usepackage{bbm}
\usepackage{tikz-cd}

\hypersetup{
    colorlinks=true,
    linkcolor=blue,
    filecolor=magenta,
    urlcolor=magenta,
    pdftitle={MAT 535 HW}
}

\usepackage{mathtools}
\usepackage{thmtools}
\usepackage[textsize=footnotesize]{todonotes}
\usepackage{titling}
\usepackage{cancel}

\reversemarginpar
\mdfdefinestyle{mdbluebox}{roundcorner=10pt,innerbottommargin=9pt,
linecolor=blue,backgroundcolor=TealBlue!5,}
\declaretheoremstyle[headfont=\sffamily\bfseries\color{MidnightBlue},
mdframed={style=mdbluebox},]{thmbluebox}

\mdfdefinestyle{mdbloobox}{frametitlefont=\bfseries,innerbottommargin=8pt,
nobreak=true,backgroundcolor=TealBlue!5,linecolor=blue,}
\declaretheoremstyle[headfont=\bfseries\color{MidnightBlue},
mdframed={style=mdbloobox},headpunct={\\[3pt]},postheadspace=0pt,]{thmbloobox}

\mdfdefinestyle{mdredbox}{frametitlefont=\bfseries,innerbottommargin=8pt,
nobreak=true,backgroundcolor=Salmon!5,linecolor=RawSienna,}
\declaretheoremstyle[headfont=\bfseries\color{RawSienna},
mdframed={style=mdredbox},headpunct={\\[3pt]},postheadspace=0pt,]{thmredbox}

\mdfdefinestyle{mdgreenbox}{linecolor=ForestGreen,backgroundcolor=ForestGreen!5,
linewidth=2pt,rightline=false,leftline=true,topline=false,bottomline=false,}
\declaretheoremstyle[headfont=\bfseries\sffamily\color{ForestGreen!70!black},
mdframed={style=mdgreenbox},headpunct={ --- },]{thmgreenbox}

\mdfdefinestyle{mdorangebox}{linecolor=RedOrange,backgroundcolor=RedOrange!5,
linewidth=2pt,rightline=false,leftline=true,topline=false,bottomline=false,}
\declaretheoremstyle[headfont=\bfseries\sffamily\color{RedOrange!70!black},
mdframed={style=mdorangebox},headpunct={ --- },]{thmorangebox}

\mdfdefinestyle{mdblackbox}{linecolor=black,backgroundcolor=RedViolet!5!gray!5,
linewidth=3pt,rightline=false,leftline=true,topline=false,bottomline=false,}
\declaretheoremstyle[mdframed={style=mdblackbox}]{thmblackbox}

\declaretheoremstyle[spaceabove=6pt,spacebelow=6pt]{boldhead}
\declaretheoremstyle[spaceabove=6pt,spacebelow=6pt,headfont=\normalfont\itshape]{italichead}

\declaretheorem[style=boldhead,name=Problem]{problem}
\declaretheorem[style=boldhead,name=Problem,numbered=no]{problem*}
\declaretheorem[style=boldhead,name=Setup]{setup} % for config geo
\declaretheorem[style=boldhead,name=Example,sibling=problem]{example}
\declaretheorem[style=boldhead,name=$\bigstar$ Problem,sibling=problem]{niceprob}
\declaretheorem[style=boldhead,name=Theorem,sibling=problem]{theorem}
\declaretheorem[style=boldhead,name=Corollary,sibling=theorem]{corollary}
\declaretheorem[style=boldhead,name=Proposition,sibling=theorem]{prop}
\declaretheorem[style=boldhead,name=Lemma,numberwithin=problem]{lemma}
\declaretheorem[style=boldhead,name=Theorem,numbered=no]{theorem*}
\declaretheorem[style=boldhead,name=Lemma,numbered=no]{lemma*}
\declaretheorem[style=boldhead,name=Claim,numberwithin=problem]{claim}
\declaretheorem[style=boldhead,name=Claim,numbered=no]{claim*}
\declaretheorem[style=boldhead,name=Remark,numberwithin=problem]{remark}
\declaretheorem[style=boldhead,name=Remark,numbered=no]{remark*}
\declaretheorem[style=boldhead,name=Definition,sibling=theorem]{definition}
\declaretheorem[style=boldhead,name=Definition,numbered=no]{definition*}

\providecommand{\ol}{\overline}
\providecommand{\tensor}{\otimes}
\renewcommand{\hom}{\operatorname{Hom}}
\providecommand{\tor}{\operatorname{Tor}}
\providecommand{\Gal}{\operatorname{Gal}}
\providecommand{\ceil}[1]{\left\lceil #1 \right\rceil}
\providecommand{\floor}[1]{\left\lfloor #1 \right\rfloor}
\newcommand{\dd}{\,\mathrm{d}}
\providecommand{\ul}{\underline}
\providecommand{\eps}{\varepsilon}
\providecommand{\half}{\frac{1}{2}}
\providecommand{\CC}{\mathbb C}
\providecommand{\FF}{\mathbb F}
\providecommand{\II}{\mathbb I}
\providecommand{\NN}{\mathbb N}
\providecommand{\QQ}{\mathbb Q}
\providecommand{\RR}{\mathbb R}
\providecommand{\ZZ}{\mathbb Z}
\providecommand{\ind}{\mathbbm 1}
\providecommand{\dg}{^\circ}
\providecommand{\ii}{\item}
\providecommand{\setdiff}{\smallsetminus}
\providecommand{\alert}[1]{{\sffamily\textbf{\textcolor{blue}{#1}}}}
\providecommand{\inv}{^{-1}}
\providecommand{\mb}[1]{\mathbf{#1}}
\newcommand{\what}{\widehat}

\newenvironment{soln}{\begin{proof}[Solution]}{\end{proof}}

\providecommand{\pow}{\operatorname{Pow}}
\providecommand{\vocab}[1]{{\textbf{\textcolor{ForestGreen}{#1}}}}
\providecommand{\lcm}{\operatorname{lcm}}
\providecommand{\abs}[1]{\left\lvert #1\right\rvert}
\providecommand{\smabs}[1]{\lvert #1\rvert}
\providecommand{\innerprod}[1]{\langle\!\langle #1\rangle\!\rangle}
\providecommand{\norm}[1]{\left\| #1\right\|}
\providecommand{\smnorm}[1]{\| #1\|}
\providecommand{\gr}{\operatorname{gr}}

\usepackage{mathrsfs}

% place title at top right
\setlength{\droptitle}{-8ex}
\pretitle{\begin{flushright}\Large\bfseries}
\posttitle{\par\end{flushright}}
\preauthor{\begin{flushright}}
\postauthor{\end{flushright}}
\predate{\begin{flushright}}
\postdate{\end{flushright}}

\title{MAT 535 HW10}
\author{\large Aditya Pahuja}
\date{\large Due 04/30}
\begin{document}
\maketitle

\begin{problem}
    [\S18.1\#7]
    Let $V$ be be $4$-dimensional permutation module for $S_4$
    described in Example $3$ of the second set of examples.
    Let $\pi\colon D_8\to S_4$ be the permutation representation of $D_8$
    obtained from the action of $D_8$ by
    left multiplication on the set of left cosets of its subgroup
    $\langle s\rangle$. Make $V$ into an $FD_8$-module via $\pi$
    as described in Example $4$ and write out the $4\times 4$ matrices
    for $r$ and $s$ given by this representation
    with respect to the basis $e_1$, $e_2$, $e_3$, $e_4$.
\end{problem}

\begin{soln}
    The left cosets of $\langle s\rangle$ are the subgroup itself,
    $\{r,rs\}$, $\{r^2,r^2s\}$, and $\{r^3,r^3s\}$,
    which we will say correspond respectively to
    the basis elements $e_1$, $e_2$, $e_3$, $e_4$.
    Then, $r$ acts by cyclically permuting the cosets, so
    the matrix corresponding to this action is
    \begin{equation*}
        \rho(r) = \begin{pmatrix}
	    0&0&0&1\\
	    1&0&0&0\\
	    0&1&0&0\\
	    0&0&1&0
        \end{pmatrix}
    \end{equation*}
    Left multipliction by $s$ preserves $\langle s\rangle$ and $\{r^2, r^2s\}$
    and swaps $\{r,rs\}$ and $\{r^3, r^3s\}$ due to the fact $sr^k = r^{-k}s$,
    so the matrix describing the action of $s$ is
    \begin{equation*}
        \rho(s) = \begin{pmatrix}
	    1&0&0&0\\
	    0&0&0&1\\
	    0&0&1&0\\
	    0&1&0&0\\
        \end{pmatrix}\qedhere
    \end{equation*}
\end{soln}

\begin{problem}
    [\S18.1\#13]
    Let $R$ be a ring and let $M$ and $N$ be simple $R$-modules.
    \begin{enumerate}[(a)]
        \ii Prove that every nonzero $R$-module homomorphism
	from $M$ to $N$ is an isomorphism.
	\ii Prove Schur's lemma: if $M$ is a simple $R$-module
	then $\hom_R(M,M)$ is a division ring.
    \end{enumerate}
\end{problem}

\begin{soln}
    (a) The kernel of a nonzero homomorphism from $M$ to $N$
    cannot be $M$, and the only other choice of submodule
    is $0$, so the homomorphism is injective.
    Similarly, the image of the nonzero homomorphism
    cannot be zero, so it must be $M$,
    so the homomorphism is surjective.

    (b) Every nonzero element of $\hom_R(M,M)$
    is an isomorphism by (a), hence is invertible.
\end{soln}

\begin{problem}
    [\S18.2\#11]
    Prove that if $R$ is a ring with $1$
    such that every $R$-module is free then $R$ is a division ring.
\end{problem}

\begin{soln}
    We wish to show that each $a\in R-\{0\}$ has a two-sided inverse.
    Let $I = Ra$ be the left ideal generated by $a$;
    by assumption, the $R$-modules $I$ and $R/I$ are both free.
    We can construct a section $R/I\to R$ of the projection $R\to R/I$
    by choosing the images of a basis of $R/I$,
    so that the short exact sequence $0\to I\to R\to R/I\to 0$ splits,
    where the map $I\to R$ is the inclusion.
    The ranks of $I$ and $R/I$ as $R$-modules thus add to
    the rank of $R$, which is $1$; since $I\ne 0$ (it contains $a$),
    $R/I = 0$, so that the inclusion $I\to R$ is an isomorphism,
    which implies $I = R$. Thus multiplication on the right by $a$
    is a permutation of $R - \{0\}$ sending an element $a\inv$ to $1$;
    then right multiplication by $a\inv$ sends $a$ to $1$ as well,
    so $a\inv$ is the (two-sided) inverse of $a$.
\end{soln}

\pagebreak

\begin{problem}
    [\S18.2\#12]
    Let $F$ be a field, let $f(x)\in F[x]$, and let $R = F[x]/(f(x))$.
    Find necessary and sufficient conditions on the factorization
    of $f(x)$ in $F[x]$ so that $R$ is a semisimple ring.
\end{problem}

\begin{soln}
    A commutative ring is semisimple
    if and only if it is a product of fields
    because a matrix ring over a division algebra
    is commutative if and only if it is a field.
    Let $f(x) = p_1(x)^{e_1}\cdots p_r(x)^{e_r}$
    for pairwise distinct irreducible polynomials
    $p_i(x)$ and positive integers $e_i$.
    By the Chinese remainder theorem,
    \begin{equation*}
	R\cong\prod_{i=1}^r F[x]/(p_i(x)^{e_i}).
    \end{equation*}
    A direct product of fields has no nonzero nilpotent elements,
    so each $e_i$ must be $1$ (else $p_1(x)\cdots p_r(x)$
    would be a nonzero nilpotent), and in this case
    the factors given by the Chinese remainder theorem
    are indeed fields, so $f(x)$ being separable
    is necessary and sufficient for $R$ to decompose as
    a product of fields.
\end{soln}

\begin{problem}
    [\S18.3\#13]
    Let $G$ be the cyclic group of order $n$ and let $R = \QQ G$.
    Describe the Wedderburn decomposition of $R$
    and find the number and degrees of the irreducible representations
    of $G$ over $\QQ$. In particular, show that if $n = p$ is a prime
    then $G$ has exactly one nontrivial irreducible representation
    over $\QQ$ and that this representation has degree $p-1$.
\end{problem}

\begin{soln}
    The group ring $\QQ G$ is isomorphic to $\QQ[x]/(x^n - 1)$.
    By the previous problem this ring is semisimple
    (since $x^n - 1$ is separable), and
    the factorization of $x^n - 1$ over $\QQ$
    is into cyclotomic polynomials, so
    \begin{equation*}
	\QQ G \cong \prod_{d\mid n} \QQ[x]/(\Phi_d(x)).
    \end{equation*}
    Thus there is one irreducible representation
    of degree $\varphi(d)$ for each divisor $d$ of $n$.
    In particular, when $n = p$ is prime,
    there is an irreducible representation of degree $\varphi(1) = 1$
    and one of degree $\varphi(p) = p-1$.
\end{soln}










\end{document}
